\scanpage{95}
\mspage{94}

formuler ainsi, en termes d'opérations sur des parties $A,B,A_i$ de $M$, pas néc. fermées, et formées d'objets mutuellement compatibles.

Commençons par d) (\unsure{O.P.S.|Q.P.S.} $A=\widetilde F$)
\[
S_{\widetilde F}\overset{\mathrm{def}}=\operatorname{Supp}\left(\bigcup_{X\in\widetilde F}\unsure{X^\circ|X^0|X^o}\right)=\operatorname{Supp}\deleted{\unsure{X}}F
\]
ce qui revient, par définition des supports, à avoir, pour $\forall Y\in M$
\diagram{\begin{tikzcd}[ampersand replacement=\&,column sep=large,row sep=large]
Y\parallel\deleted{\unsure{X}}F \arrow[r,Leftrightarrow] \arrow[d,Leftrightarrow] \& \forall X\in\widetilde F,\ Y\parallel X^\circ \\
(Y\parallel\widetilde F) \&
\end{tikzcd}}
et cela résulte aussitôt du cor. au lemme 1.

Voyons a) i.e
\[
S_A\subset S_B\quad\Rightarrow\quad A\subset B
\]
Soit donc $X\in A$, on doit prouver $X\in B$, on est ramené à prouver que
\diagram{\begin{tikzcd}[ampersand replacement=\&,column sep=small,row sep=small]
S_{\{X\}} \arrow[d,equal] \& {}\subset S_B\quad\Rightarrow\quad X\in B \\
\operatorname{supp}(X^\circ) \&
\end{tikzcd}}
i.e.

\uline{Lemme \unsure{3|8}} Soient $X\in M$, $B\subset M$, tels que $B\cup\{X\}$ soit admissible (formé d'objets mutuellement admissibles) alors
\[
X^\circ\in S_B\quad\Rightarrow\quad X\in B
\]
% slanted
\margin{\uline{NB} Cela \deleted{\unsure{implique|montre}} comment \unsure{récupérer} $B$ en termes de $S_B$}
Le premier mb \deleted{\unsure{implique}, \unsure{en prenant}} comporte
\[
\forall Z\ \deleted{\unsure{\in M|\in A}}\text{ tel que }Z\parallel\bigcup_{Y\in B}Y^\circ,\quad\text{on a }Z\parallel\deleted{\unsure{Y}}X^\circ
\]
